\title{Byte Latent Transformer: Patches Scale Better Than Tokens}

\begin{abstract}
We present the Byte Latent Transformer ({\textsc BLT}), a novel architecture for byte-level large language models that achieves, for the first time, performance at scale comparable to token-based LLMs, while offering notable gains in inference robustness and computational efficiency. {\textsc BLT} transforms input bytes into variable-length patches that constitute the core computational elements of the model. Patch segmentation is determined by analyzing the entropy of forthcoming bytes, allowing the model to assign greater computational resources and capacity in regions of higher informational complexity. We conduct the first scaling analysis for byte-level models under controlled floating point operation constraints, extending up to 8B parameters and 4T bytes of training data. Our findings establish the practicality of scaling language models directly from raw bytes without reliance on a pre-defined vocabulary. The approach increases both training and inference efficiency through adaptive use of longer patches in more predictable data segments, resulting in qualitative advances in both reasoning capabilities and generalization to rare cases. In summary, at equivalent inference cost, {\textsc BLT} achieves superior scaling characteristics compared to tokenization-based models by jointly expanding both the patch size and the underlying model capacity.
\end{abstract}

\section{Introduction}
\label{section:intro}

We present the Byte Latent Transformer~(\textbf{BLT}), an architecture that forgoes tokenization by directly ingesting sequences of raw byte data. BLT achieves parity with token-based models at scale in terms of performance, while offering marked gains in efficiency and robustness (\S\ref{sec:robustness}), a capability not previously demonstrated. In prevailing large language models (llms), nearly all training is performed in an end-to-end manner except for the step of tokenization—a heuristic that partitions byte sequences into a fixed token inventory. This token-based approach influences the representation and compression of strings, introducing several weaknesses: sensitivity to domain or modality shifts~\citep{dagangetting}, increased vulnerability to input perturbations~(\S\ref{sec:robustness}), limited handling of orthographic features~\citep{edman2024cute}, and persistent inequities across languages~\citep{liang2023xlm,petrov2024language,limisiewicz2024myte}.

Historically, tokenization has played a vital role because training llms directly on raw bytes is computationally infeasible at scale, owing to the resulting lengthy input sequences~\citep{xue2022byt5}. Previous studies have addressed this issue by adopting more efficient self-attention mechanisms~\citep{el2020characterbert, clark2022canine} or architectures that eliminate attention altogether~\citep{wang2024mambabyte}~(\S\ref{section:related_work}). Nevertheless, these strategies primarily benefit the training of \textit{smaller models}. In large-scale settings, the majority of a Transformer's computational burden stems from the expansive feed-forward network layers operating across every byte, whereas the attention mechanism's cost becomes relatively less significant.

To optimize the use of computational resources, we introduce a dynamic, trainable strategy for aggregating bytes into \textit{patches}~(\S\ref{section:patching}), accompanied by a novel model design that integrates both byte-level and patch-level information. In contrast to traditional tokenization approaches, {\textsc BLT} does not employ a predetermined patch vocabulary. Instead, arbitrary byte groupings are transformed into latent patch embeddings using compact, parameterized encoder and decoder modules. Our results demonstrate that this approach allows for a \textit{more} efficient use of compute compared to models relying on tokenization.

Tokenization-based language models assign identical computational resources to each token. This prioritization favors uniformity over efficiency, as the tokenization process employs compression strategies that do not consistently align with the intricacy of prediction tasks. At the core of our proposed architecture lies the principle that computational effort should be distributed adaptively, according to the specific demands of the input. For instance, utilizing a sizable transformer is unnecessary for predicting the final characters in most words, as these represent relatively simple and predictable—i.e., low-entropy—tasks in comparison to selecting the initial word of a new sentence. This concept is embodied in the architecture of {\textsc BLT}~(\S\ref{section:architecture}), which utilizes three transformer blocks: two lightweight byte-level \textit{local models}, alongside a more substantial global \textit{latent transformer}~(\autoref{fig:arch_high_level}). 

In order to organize bytes into patches and adjust computation dynamically, {\textsc BLT} segments data by analyzing the entropy associated with next-byte prediction. This method produces contextually informed groupings of bytes characterized by comparable information content.

We conduct the first study that scales byte-level models under a controlled flop budget, extending to models of up to 8B parameters and utilizing 4T training bytes. Our results demonstrate the feasibility of large-scale, end-to-end training on bytes, eliminating the necessity for predetermined vocabulary tokenization. In terms of training performance per flop, {\textsc BLT} achieves parity with Llama 3, yet requires up to 50\% less computational cost at inference time (\S\ref{section:scaling}).\footnote{Model computational requirements are quantified as the total number of Floating Point OPerations (flops) executed.} 

Furthermore, operating directly on raw bytes leads to notable gains in handling infrequent data patterns. {\textsc BLT} exhibits increased resilience to noisy data compared to models that rely on tokenizers, along with superior abilities at interpreting characters, as substantiated by tests in orthographic processing, phonological tasks, and low-resource machine translation (\S\ref{sec:robustness}). 

Moreover, with {\textsc BLT}, one can expand both the model and input patch sizes simultaneously while keeping the inference flop budget unchanged. Employing larger patch sizes typically reduces computation required during inference, thereby freeing up resources that allow for scaling up the global latent transformer, which operates less frequently. Our inference-flop constrained scaling evaluations (\autoref{fig:fixed_inference_scaling}) reveal that such byte-level, non-tokenized models scale more favorably than those employing conventional tokenization mechanisms.

To conclude, the contributions of this work are as follows: 1) We present {\textsc BLT}, a byte latent llm framework that enables adaptive compute allocation to enhance flop efficiency, 2) We demonstrate parity with Llama 3 in terms of training flop expenditure at the 8B parameter scale, while permitting a flexible trade-off that achieves up to 50\% greater flop efficiency with only modest impacts on downstream evaluation metrics, 3) {\textsc BLT} introduces a novel axis for llm scaling, allowing model capacity to increase without exceeding a fixed inference compute budget, and 4) We provide evidence that {\textsc BLT} models exhibit stronger resilience to input noise and better sensitivity to sub-word characteristics in input data, outperforming token-based llms in these respects. All code for training and inference with {\textsc BLT} is publicly available at~\url{https://github.com/facebookresearch/blt}.

\section{Patching: From Individual Bytes to Groups of Bytes}
\label{section:patching}

Dividing bytes into \textit{patches} enables {\textsc BLT} to adjust computational resources flexibly, depending on the surrounding context. Figure~\ref{fig:patching_types} illustrates various strategies for grouping bytes into patches. More formally, a patching function $f_p$ partitions a byte sequence $\pmb{x}=\{x_i,|i=1,\ldots n\}$ of length $n$ into $m < n$ patches $\pmb{p}=\{p_j|j=1,\ldots,m\}$ by assigning each $x_i$ to the set \{0,1\}, where a value of 1 marks the initiation of a new patch. In both token-level and patch-level frameworks, the main Transformer’s step count—and thus the bulk of computational expense—relates directly to the number of segments formed, i.e., patches in the case of {\textsc BLT} given a particular patching strategy. Thus, the expected length of a patch, referred to as the \textit{patch size}, becomes the key determinant of the processing cost during both the training phase and inference for a specific patching function~(\S\ref{section:flops}).

We proceed by describing three specific patching approaches: fixed-length patching by grouping a constant number of bytes per patch~(\S\ref{section:static-patch}), patching based on whitespace characters~(\S\ref{section:white-patch}), and adaptive patching leveraging byte-level language model entropies~(\S\ref{section:dyn-patch}). In conclusion, we review incremental patching and clarify the distinction between tokenization and patching~(\S\ref{section:bpe-patch}).

\subsection{Strided Patching Every K Bytes}
\label{section:static-patch}

A common technique for organizing bytes is to segment them into fixed-size patches of $k$ bytes, as implemented in MegaByte~\citep{yu2023megabyte}. This uniform partitioning allows for straightforward implementation during both training and inference, simplifies adjustment of average patch size, and thereby facilitates easy management of computational requirements. Despite these advantages, this approach has notable limitations. One major issue is the lack of adaptive computation: computational resources are not allocated based on local content demands—some transformer steps $j$ may be wasted on uninformative regions like code whitespace, while insufficient compute might be allocated to content-rich bytes such as mathematical expressions. Additionally, this approach may cause inconsistent and non-context-aware segmentation of identical byte sequences; for example, the same word may be divided differently across patches.

\subsection{Space Patching}
\label{section:white-patch}

\citet{slagle2024spacebyte} introduces a straightforward enhancement to strided patching by generating new patches at every occurrence of a space-like byte\footnote{
    Space-like bytes are those that are neither latin characters, digits, nor utf-8 continuation bytes; additionally, each resulting patch must include at least one non space-like byte.}, which frequently serve as the boundaries between linguistic units in various languages. In the Space patching technique, an additional latent transformer computation (corresponding to increased FLOPs) is dedicated to modeling each word individually. This approach guarantees a consistent patching of words across sequences and strategically allocates more computational resources to challenging prediction points, which are typically found immediately after space characters. For instance, identifying the initial byte of the answer in “Who composed the Magic Flute?\uline{\hspace{.6em}}” poses a much greater challenge than predicting subsequent bytes once the leading "M" has been established, since the initial character narrows the field of plausible continuations dramatically—subsequent characters in “Mozart” are thus easier to forecast. Nonetheless, space patching presents limitations: it does not adapt well to all languages or content domains and lacks the ability to vary patch sizes. In the following, we present a novel patching approach inspired by the observation that initial bytes within words tend to be most predictive, while also enabling dynamic control of patch size.

\subsection{Entropy Patching: Using Next-Byte Entropies from a Small Byte LM}
\label{section:dyn-patch}

Instead of depending on rule-based heuristics like whitespace, we propose a data-driven method to detect next-byte predictions characterized by high uncertainty. Specifically, we present \textit{entropy patching}, a technique that establishes patch boundaries utilizing entropy estimates.

We train a small byte-level auto-regressive language model on the training data for {\textsc BLT} and compute next byte entropies under the LM distribution $p_{e}$ over the byte vocabulary $\mathcal{V}$:

\begin{equation}
H(x_i) = \sum_{v \in \mathcal{V}} p_{e}(x_i=v|\pmb{x}_{<i}) \log p_{e}(x_i=v|\pmb{x}_{<i})
\end{equation}

We utilize two approaches for detecting patch boundaries based on the entropies $H(x_i)$. The initial method involves selecting points that surpass a fixed global entropy threshold, as depicted in~\autoref{fig:patching}. The alternative approach focuses on finding points where the entropy is comparatively higher than the preceding value. This latter method can also be viewed as recognizing instances where the approximate monotonic decrease in entropy within a patch is disrupted.

\begin{align*}
    \text{Global Constraint}\hspace{1cm}H(x_t)& > \theta_g\\
    \text{Approx. Monotonic Constraint}\hspace{1cm}H(x_t)& - H(x_{t-1}) > \theta_r
\end{align*}

Patch boundaries are determined through a lightweight preprocessing phase that occurs while data is being loaded. This approach contrasts with~\citet{nawrot-etal-2023-efficient}, where a classifier is employed to detect entropy-driven patch boundaries. Within our experiments~(\S\ref{section:experiments}), we evaluate and contrast these two techniques for classifying bytes with low versus high entropy.

\subsection{The Byte-Pair Encoding (BPE) Tokenizer and Incremental Patching}
\label{section:incremental-patching}
\label{section:bpe-patch}

A large number of contemporary llms, such as our baseline Llama 3, employ subword tokenization methods like bpe~\citep{gage1994new, sennrich-etal-2016-neural}. In our terminology, ``tokens'' denote fixed byte-group units chosen from a \textit{finite} vocabulary established before training, distinguishing them from ``patches,'' which are created dynamically from sequences and do not rely on a predetermined vocabulary. Notably, a significant distinction is that when using tokens, the underlying byte-level information is not directly accessible to the model.

An important advancement that {\textsc BLT} offers compared to tokenization-based approaches is the way it reshapes the balance between vocabulary size and computational requirements. In conventional LLMs, expanding the vocabulary size leads to tokens that encode more information on average, thus reducing the number of decoding steps; however, this also results in a larger output space for the model’s final projection layer. This balance constrains the extent to which tokenization-based methods can alter both token granularity and the associated inference expenditure. To illustrate, Llama 3 achieves an increase in average token length from 3.7 to 4.4 bytes, but this improvement comes with a fourfold enlargement of its embedding matrix when compared to Llama 2.

During generation, {\textsc BLT} must determine whether the present point in the byte sequence corresponds to a patch boundary, as this choice dictates if additional computation through the Latent Transformer is required. Importantly, this determination must be made without knowledge of the remainder of the sequence that has not yet been produced. Consequently, the process of patching must avoid relying on any subsequent bytes to decide how to segment the sequence. Mathematically, a patching scheme $f_p$ is said to possess the property of incremental patching if the following holds:
$$f_p(\pmb{x}_{<i}) = f_p(\pmb{x})_{<i}$$
The bpe method does not qualify as an incremental patching scheme since the same initial segment can receive different tokenizations depending on what follows it, thus failing to meet the above requirement\footnote{Introducing a dedicated delimiter token to mark patch boundaries can modify bpe into an incremental patching scheme; however, this increases the length of the resulting byte sequence.}.

\section{{\textsc BLT} Architecture}
\label{section:architecture}
{\textsc BLT} consists of a sizable global autoregressive language model responsible for processing patch representations, paired with two more compact local modules. These local components handle the encoding of byte sequences into patches and the decoding of patch representations back to bytes~(\autoref{fig:arch_high_level}).

\subsection{Latent Global Transformer Model}

\textit{The Latent Global Transformer} refers to an autoregressive transformer architecture, denoted as $\mathcal{G}$, composed of $l_{\mathcal{G}}$ layers. This model takes in a series of latent input patch representations, $p_j$, and transforms them into corresponding output patch representations, $o_j$. In this work, $j$ is consistently used as the index for patches, while $i$ indexes bytes. The global transformer applies a block-causal attention mask~\citep{dubey2024llama}, limiting each patch’s attention to its predecessors and itself within the same document. The majority of the computational workload in both pre-training and inference is due to this component; therefore, controlling the frequency with which it is applied offers a means to adapt computational resource allocation in relation to the input and output complexity across different segments.

\subsection{Local Encoder}
\label{section:local}

The \textit{Local Encoder Model}, denoted as $\mathcal{E}$, comprises a compact, transformer-based architecture characterized by a significantly smaller number of layers ($l_{\mathcal{E}} << l_{\mathcal{G}}$) compared to the global counterpart. Its principal function is to rapidly convert a sequence of input bytes $b_i$ into rich patch-level representations $p_j$. Distinguishing itself from standard transformers, this model incorporates a cross-attention layer subsequent to each transformer layer to aggregate byte-level embeddings into patch-level vectors (see~\autoref{fig:crossattn}).
Initially, the byte sequence $b_i$ is projected via an embedding matrix of size $\mathbb{R}^{256 \times h_{\mathcal{E}}}$, resulting in representations $x_i$. These initial embeddings can be optionally enhanced with supplementary information through hash-embeddings~(\S\ref{sec:hash-emb}). 
Subsequently, the model applies alternating layers of transformer blocks and cross-attention modules~(\S\ref{sec:enc-cross-attn}) to generate patch representations, $p_i$, which are subsequently utilized by the global transformer, $\mathcal{G}$. The attention mechanism employed in the transformer layers follows a \textit{local block causal} scheme, in which each byte can attend to a fixed number $w_{\mathcal{E}}$ of preceding bytes; this window can extend across evolving patch boundaries but does not span document limits. The following sections further elaborate on the design of the embedding process and the cross-attention module.

\subsubsection{Encoder Hash n-gram Embeddings}
\label{sec:ngram-emb}
\label{sec:hash-emb}
An essential aspect of forming rich and resilient representations at every position $i$ is to embed contextual knowledge about prior bytes. Within {\textsc BLT}, this is addressed by capturing characteristics of the current byte $b_i$ on its own \textit{as well as} within the context of a byte n-gram. Specifically, at each time step $i$, we initially generate corresponding byte-grams

\begin{align}
    g_{i,n}=\{b_{i-n + 1},\ldots, b_i\}
\end{align}

for every byte position $i$ and for $n$ ranging from three to eight.\footnote{
Byte-grams are not considered for sizes $n$ or greater whenever $i<n$. }

Next, we propose hash-based $n$-gram embeddings, assigning all byte $n$-grams to indices in a fixed-size embedding table $E_{n}^{hash}$, with each $n$ in $\{3,4,5,6,7,8\}$, determined through a hash function~\citep{bai2010learning}. The computed embedding for a given $n$-gram is then added to the byte’s embedding. This sum is normalized prior to serving as input to the local encoder. The enhanced embedding is thus computed as follows:

\begin{align}
e_i &= x_i + \sum_{n = 3,...,8}E_{n}^{hash}(\text{Hash}(g_{i,n})) \\
\text{where, Hash}(g_{i,n}) &= \text{RollPolyHash}(g_{i,n}) \% |E_{n}^{hash}|
\end{align}

We scale $e_i$ by dividing by the number of $n$-gram sizes incremented by one, applying RollPolyHash as outlined in Appendix~\ref{appendix:rollpolyhash}. Section~\ref{section:ablations} investigates how varying the values of $n$ and the size of the embedding table for $n$-gram hash embeddings influence the scaling behavior under flop constraints. Alongside hash-based $n$-gram embeddings, we also conducted experiments employing frequency-based $n$-gram embeddings; further information regarding this line of inquiry is presented in Appendix~\ref{appendix:ngrams}.

\subsubsection{Encoder Multi-Headed Cross-Attention}
\label{sec:enc-cross-attn}
Our design largely replicates the input cross-attention mechanism described in the Perceiver architecture~\citep{jaegle2021perceiver}, with the principal modification being that latent variables are aligned with the representations of variable-sized patches rather than a static group of latent vectors~(\autoref{fig:crossattn}), and each latent vector is restricted to attending solely to the bytes contained in its associated patch. Within the module, each patch $p_j$ has an assigned query vector, which is formed by aggregating the byte representations within $p_j$ and subsequently applying a linear transformation, $\mathcal{E}_{C} \in \mathbb{R}^{h_{\mathcal{E}} \times (h_{\mathcal{E}} \times U_{\mathcal{E}})}$, where $U_{\mathcal{E}}$ specifies the count of encoder cross-attention heads. To express this formally, suppose $f_{\text{bytes}}(p_j)$ provides the sequence of bytes for patch $p_j$, then we obtain

\begin{align}
P_{0,j} &= \mathcal{E}_{C}(f_\text{bytes}((p_j)), f~ \text{is a pooling function} \\
P_l &= P_{l-1} + W_o\left(\text{softmax}\left(\frac{QK^T}{\sqrt{d_k}}\right)V\right) \\
\text{where } Q_j &= W_q(P_{l-1,j}), K_i = W_k(h_{l-1,i}), V_i = W_v(h_{l-1,i}) \\
h_l &= \text{Encoder-Transformer-Layer}_l(h_{l-1})
\end{align}

Here, $P \in \mathbb{R}^{n_p \times h_{\mathcal{G}}}$ denotes $n_p$ patch-level representations intended for processing by the global model, with the initial values obtained by aggregating the byte embeddings $e_i$ that belong to each patch $p_j$. The matrices $W_q$, $W_k$, $W_v$, and $W_o$ serve as projection parameters for queries, keys, values, and outputs, where the keys and values are generated by applying the projections to the byte-level representations $h_i$ derived from the preceding layer (or directly from $e_i$ when considering the first layer). Our masking approach aligns specifically with the patch structure, ensuring that each query vector $Q_j$ can only attend to those key and value vectors that originate from the bytes within the same patch $j$. Given that we deploy multi-head attention over $Q, K$, and $V$, and patch representations frequently possess higher dimensionality ($h_{\mathcal{G}}$) than $h_{\mathcal{E}}$, we retain $P_l$ as multiple attention heads each with dimension $h_{\mathcal{E}}$ for the cross-attention operation; these are subsequently concatenated to achieve the full $h_{\mathcal{G}}$ dimensionality. Moreover, queries, keys, and values undergo pre-LayerNorm normalization, and positional embeddings are omitted within this cross-attention mechanism. A residual connection wraps around the entire cross-attention module.

\subsection{Local Decoder} 

Analogous to the design of the local encoder, the local decoder $\mathcal{D}$ utilizes a compact transformer architecture featuring $l_{\mathcal{D}} << l_{\mathcal{G}}$ layers. Its purpose is to reconstruct the sequence of raw bytes, $y_i$, from the global patch representations, $o_j$. The decoding process generates the byte sequence sequentially, conditioning each prediction on the previously decoded bytes. For this task, the local decoder employs the hidden states outputted by the local encoder for the respective byte sequence. The overall architecture consists of $l_{\mathcal{D}}$ layers that alternate between cross-attention and transformer modules. Within each layer, cross-attention precedes the transformer operation, transforming patch representations into preliminary byte embeddings. Subsequently, the transformer layer refines these embeddings over the growing byte sequence.

\subsubsection{Decoder Multi-headed Cross-Attention} 

In the decoder cross-attention, the roles of the queries and key/values are interchanged i.e. the byte-representations are now the queries, and the patch representations are now the key/values. The initial byte-representations for the cross-attention are initialized as the byte embeddings from the last encoder layer i.e. $h_{l_{\mathcal{E}}}$. The subsequent byte-representations for layer $l$, $d_{l,i}$ are computed as:

\begin{align}
D_0 &= h_{l_{\mathcal{E}}} \\
B_l &= D_{l-1} + W_o\left(\text{softmax}\left(\frac{QK^T}{\sqrt{d_k}}\right)V\right), \\
  &\text{where } Q_i = W_q(d_{l-1,i}), K_i = W_k(\mathcal{D}_C(o_j)), V_i = W_v(\mathcal{D}_C(o_j)) \\
  D_l &= \text{Decoder-Transformer-layer}_l(B_l)
\end{align}

In this formulation, $W_k$ and $W_v$ denote the projection matrices for keys and values, respectively; these matrices perform linear projections followed by the split operation $\mathcal{D}_C$ on the terminal patch representations $o_j$ produced by the global model. The projection matrix $W_q$ is applied to the byte-level inputs $d_{l-1}$ from the preceding decoder transformer layer (or $h_{l_{\mathcal{E}}}$ in the case of the initial layer), serving as the query transformation. $W_o$ is designated as the output projection matrix. Consequently, $B \in \mathbb{R}^{h_{\mathcal{D}} \times n_b}$, where $n_b$ signifies the total count of output bytes. Subsequently, the representations in the following decoder layer, $D_l$, are derived by applying a decoder transformer layer to $B$, the output from the cross-attention mechanism. This architecture mirrors the cross-attention employed in the local encoder: it utilizes multiple attention heads, adopts pre-layer normalization, omits positional embeddings, and introduces a residual connection encasing the cross-attention component.

\section{Experimental Setup}
\label{section:experiments}
We conduct a series of rigorously controlled experiments to contrast {\textsc BLT} with tokenization-based modeling approaches, making particular efforts to ensure that {\textsc BLT} receives no benefit from access to potentially larger contextual spans.

\subsection{Pre-training Datasets} The models evaluated in this study are pre-trained utilizing two primary datasets: 1) The Llama 2 dataset~\citep{touvron2023llama}, containing 2 trillion tokens sourced from an array of publicly accessible repositories, and subjected to comprehensive cleaning and filtering procedures to enhance overall data quality; and 2) {\textsc BLT}-1T: An original collection of 1 trillion tokens, assembled from diverse public sources, incorporating a portion of pre-training material published by Datacomp-LM~\citep{li2024datacomp}. The Llama 2 dataset serves in scaling law analyses, following the approach of~\cite{dubey2024llama} to guide the selection of optimal architectures for {\textsc BLT}. In contrast, the {\textsc BLT}-1T dataset supports a full-scale pre-training run, providing a benchmark to directly evaluate against Llama 3 across a range of downstream tasks. Importantly, neither dataset contains any information originating from Meta platforms or services. Additionally, in all baseline studies involving tokenizer-based models, the Llama 3 tokenizer is adopted—featuring a 128K token vocabulary—since it consistently yielded superior performance compared to the Llama 2 tokenizer under our experimental settings.

\subsection{Entropy Model} For our experiments, the entropy model is implemented as a byte-level language model, trained on the identical data distribution used for the full {\textsc BLT} model. By default, the architecture employed is a transformer consisting of 100 million parameters across 14 layers, with a hidden size of 512, utilizing sliding window attention over 512 bytes. The other hyperparameter settings match those used for both our local and global transformer models. As outlined in \autoref{section:ablations}, we evaluated a variety of model scales, receptive field sizes, and architectural variants. Notably, in cases where the receptive field is sufficiently constrained, it is possible to efficiently represent the trained entropy model within a lookup table.

\subsection{Entropy Threshold and Equalizing Context Length} For models utilizing entropy-based patching, we determine a patching threshold that results in a specified mean \textit{patch size} across the pretraining data mixture. In the case of {\textsc BLT}, unlike traditional tokenization, the \textit{patch size} is not constrained and can be set arbitrarily, which can notably impact the context window available to the model. To ensure a fair comparison and prevent models with larger patch sizes from benefiting from a longer context, we standardize the total number of bytes per batch in expectation. Consequently, for models with increased patch size, the sequence length is shortened. On the Llama 2 dataset, we configure the context window to 8k bytes; for the {\textsc BLT}-1T dataset, the context window is expanded to 16k bytes on average, while the expected batch size remains fixed at 16M bytes.

Although the average batch size remains fixed, dynamic patching techniques result in varying proportions of bytes to patches during batch loading. To optimize performance, our {\textsc BLT} training procedure organizes batches by grouping patches together, thereby eliminating the need for padding operations in the computationally intensive latent transformer. This strategy guarantees a uniform patch count per batch. In the training phase, byte sequences are padded—or, if necessary, truncated—to 12k for Llama 2 and 24k for {\textsc BLT}-1T datasets. This approach helps prevent unexpected memory surges caused by batches with exceptionally large patches.

\subsection{Entropy Model Context}
\label{section:entropy-context}
Our empirical observations indicate that entropy patching leads to increasingly extensive patches within highly structured tasks, such as multiple choice problems, which tend to exhibit significant repetition (for example, see the patching process in \autoref{fig:mmlu_patching}). These differences stem from decreased entropy in segments of the model context where repetitive patterns are prevalent. Accordingly, in the comprehensive experiment with {\textsc BLT}-Entropy and a patch size of 4.5, we periodically refresh the entropy context by introducing new lines, and employ an approximate monotonicity constraint. This approach is less susceptible to "entropy drift," which can result from variations in the length of the context. Notably, this adjustment only alters the manner in which entropies are computed; the method for determining the entropy threshold remains unchanged.

\subsection{FLOPs Estimation} 
\label{section:flops}

Our methodology for calculating transformer flops closely mirrors that of Chinchilla~\citep{hoffmann2022training}, incorporating the flops arising from feed-forward network computations, qkvo projections in self-attention mechanisms, as well as the attention calculation and output projection steps. However, one key distinction in our approach is the treatment of the input embedding layer: we consider it as an efficient lookup operation rather than a dense matrix multiplication, and thus assign it a zero flop count. Consistent with prior research, we approximate that the backwards pass incurs double the flops required for the forward pass.

To compute flops \textit{per byte} for {\textsc BLT} models, we add up the flops for the local encoder transformer, the global latent transformer, and the local decoder transformer, together with the cross attention blocks in the encoder and the decoder:

\begin{align}
    \text{FL}_{\text{{\textsc BLT}}} &= \text{Transf. FL}(h_{\mathcal{G}}, l_{\mathcal{G}}, m=n_{ctx}/n_p, V=0) / n_p\\
    &\quad + \text{Transf. FL}(h_{\mathcal{E}}, l_{\mathcal{E}}, m=w_{\mathcal{E}}, V=0) \\
    &\quad + \text{Transf. FL}(h_{\mathcal{D}}, l_{\mathcal{D}}, m=w_{\mathcal{D}}, V=256) \\ 
    &\quad + \text{Cross Attn. FL}(h_{\mathcal{E}}, l_{\mathcal{E}}, m=n_p, r=n_p/k) \cdot \frac{k}{n_p} \\ 
    &\quad + \text{Cross Attn. FL}(h_{\mathcal{D}}, l_{\mathcal{D}}, m=k, r=k/n_p)
\end{align}

Here, $n_{ctx}$ denotes the number of bytes in the input sequence, $n_p$ refers to the size of each patch, $r$ represents the proportion of queries relative to key/value pairs, and $k$ indicates the ratio between patch dimensions and byte dimensions—specifically, the count of local model partitions joined together to yield the global model representation (for example, $k=2$ as illustrated in \autoref{fig:crossattn}). The term $V$ stands for the size of the vocabulary relevant to the output projection, which is solely involved in the local decoding stage. The chosen context length $m$ for the attention mechanism varies depending on whether the module processes the byte sequence or the patch sequence. We adjust the floating-point operation (FLOP) counts of attention to reflect these distinctions for each model part. The specific formulas used to calculate Transformer-FLOPs and Cross-Attention FLOPs can be found in Appendix~\ref{section:flops-appendix}.

\subsection{Bits-Per-Byte Estimation} Perplexity only makes sense in the context of a fixed tokenizer as it is a measure of the uncertainty for each token. When comparing byte and token-level models, following previous work~\citep{xue2022byt5, yu2023megabyte, wang2024mambabyte}, we instead report Bits-Per-Byte (BPB), a tokenizer independent version of perplexity. Specifically:

\begin{align}
    \text{BPB}(x)&= \frac{\mathcal{L}_{CE}(\pmb{x})}{\ln(2)\cdot n_{\text{bytes}}}
\end{align}

Here, the overall cross-entropy loss computed on the data $\pmb{x}$ is divided by both a normalization factor and the total byte count of $\pmb{x}$, thereby quantifying the data uncertainty per byte.

\subsection{Transformer Architecture Hyperparameters} For each transformer block within {\textsc BLT}—encompassing both local and global components—we primarily adhere to the Llama~3~\citep{dubey2024llama} configuration. Specifically, the feed-forward layers utilize the SwiGLU activation function~\citep{swiglu}, self-attention mechanisms are equipped with rotary positional embeddings (RoPE)~\citep{RoPe} parameterized with $\theta=500000$~\citep{xiong2024effective}, and we employ RMSNorm~\citep{RMSNorm} as our normalization strategy. Flash Attention~\citep{dao2022flashattention} is incorporated into all self-attention layers that operate with static attention masks, including both \textit{block causal} and \textit{fixed-window block causal} variants, where the window size is fixed at 512 tokens. For cross-attention layers, due to the presence of adaptively generated, patch-specific masks, we employ Flex Attention\footnote{\url{https://pytorch.org/blog/flexattention}} for efficient, fused execution—substantially accelerating the training process.

\subsection{{\textsc BLT}-Specific Hyperparameters} 

To evaluate the performance of the {\textsc BLT} models, we pursue experiments in two main directions: analyzing scaling behavior and assessing downstream task performance. Across these investigations, we test models at varying parameter sizes: 400M, 1B, 2B, 4B, and 8B. Detailed architectural hyperparameters for each model variant are reported in Appendix~Table~\ref{tab:arch}.

For the initialization of queries in the first cross-attention stage within the local encoder, we apply max-pooling. Across all {\textsc BLT} variants, we consistently use $500,000$ hashes, employ a single hash function, and set n-gram sizes between 3 and 8. The learning rate is uniformly set to $4e-4$ across all tested model sizes. This selection was informed by a hyperparameter sweep on the 400M and 1B configurations, with rates between $1e-3$ and $1e-4$, revealing that the same learning rate is optimal for both standard token-based and {\textsc BLT} models. For scaling studies utilizing Llama-2 data, batch sizes (measured as token or byte count equivalents) are adopted as recommended by~\cite{dubey2024llama}. The models are trained with the AdamW optimizer~\citep{adamw}, configuring $\beta_1$ as 0.9, $\beta_2$ as 0.95, and using $\epsilon = 10^{-8}$. We adopt a learning rate schedule composed of a linear warm-up over 2,000 steps, followed by cosine decay to zero; a weight decay coefficient of 0.1 is applied, and global gradient norms are clipped at 1.0.

\section{Scaling Trends}
\label{section:scaling}

We provide a comprehensive overview of scaling behaviors in byte-level models that can guide future scaling of {\textsc BLT} models. Our approach to scaling seeks to overcome shortcomings in previous investigations of byte-level architectures through the following steps: (a) assessing scaling patterns within a compute-optimal training framework, (b) training 8B-parameter models using substantial datasets (reaching up to 1T tokens or 4T bytes) and evaluating their performance on a suite of downstream benchmarks, and (c) examining scaling dynamics when holding inference costs constant. Later, we will explore specific benefits that arise from modeling raw byte-sequences.

\subsection{Parameter Matched Compute Optimal Scaling Trends}
\label{section:parameter-matched-scaling}
We utilize the Llama 2 dataset to train a series of \textit{compute-optimal} bpe and {\textsc BLT} models at four distinct scales, ranging from 1B to 8B parameters. Subsequently, we chart the relationship between the total training flops and the resulting language modeling accuracy on a representative portion of the combined training datasets. The bpe models are configured according to the ideal proportion of model parameters to training examples, as prescribed by Llama~3~\citep{dubey2024llama}. This approach, rooted in the concept of \textit{compute-optimality}, is theoretically optimized to maximize performance on the training split within a predefined compute budget~\citep{hoffmann2022training}, serving as a robust baseline for comparison. For every bpe model trained, a matching {\textsc BLT} model is also trained on the identical data, employing a Latent Transformer that mirrors both the scale and design of the associated bpe Transformer.

As shown in~\autoref{fig:scaling-compute-opt} (right), {\textsc BLT} models perform on par with or better than their bpe equivalents, with this pattern persisting across varying model sizes and computational budgets. To our knowledge, {\textsc BLT} represents the inaugural byte-level Transformer framework to exhibit scaling dynamics comparable to BPE-based models under compute-optimal conditions. This finding supports our hypothesis that the ideal parameter-to-compute ratio identified for bpe models is similarly effective for {\textsc BLT}, or at the very least, does not diverge significantly.

Achieving alignment with bpe scaling trends necessitates both enhancements in architecture and the integration of dynamic patching. In ~\autoref{fig:scaling-compute-opt} (left), we present a comparison between models utilizing space-patching and Llama 3. Our approximation of SpaceByte \citep{slagle2024spacebyte} is implemented through the {\textsc BLT} space-patching approach, deliberately omitting n-gram embeddings and cross-attention mechanisms. While SpaceByte demonstrates advances over Megabyte, its performance still lags behind Llama 3. Turning to \autoref{fig:scaling-compute-opt} (right), we depict the performance gains attributable to both architectural innovations and the adoption of dynamic patching. At scale, models based on {\textsc BLT} achieve results competitive with leading tokenizer-based architectures like Llama 3.

Additionally, we find that the selection of tokenizer significantly influences performance in tokenizer-based models; specifically, models utilizing the Llama-3 tokenizer exhibit superior results compared to those employing the Llama-2 tokenizer when both are trained on identical datasets.

In summary, our {\textsc BLT} model demonstrates an architectural trend situated between Llama 2 and Llama 3 when leveraging much larger patch sizes. While the average token sizes for the bpe tokenizers in Llama 2 and Llama 3 are 3.7 and 4.4 bytes, respectively, {\textsc BLT} attains similar scaling patterns when employing mean patch sizes of 6 or even 8 bytes. Since the number of inference flops is inversely related to average patch size, selecting a patch size of 8 bytes results in an approximate reduction of inference flops by 50\%. Furthermore, increasing patch size appears to enhance performance as model scale and dataset size expand. Notably, {\textsc BLT} initialized with a patch size of 8 performs substantially worse than bpe Llama 2 at the 1B parameter level but surpasses it by the 7B scale. This pattern indicates that greater patch sizes could further improve outcomes at even larger model scales, and that scaling patch sizes alongside model capacity and compute may be advantageous.

\subsection{Beyond Compute Optimal Task Evaluations}
\label{section:evals}
To gain deeper insights into scaling behaviors, we extend the training of an 8B {\textsc BLT} model beyond the regime suggested by the compute-optimal scaling law, using the {\textsc BLT}-1T dataset—a significantly larger and higher-quality corpus. Performance is then measured on a broad array of established benchmarks, covering both classification and generation tasks. Our chosen benchmarks span the domains of common sense reasoning, general world knowledge, and code synthesis:
\paragraph{Classification tasks} comprise ARC-Easy (0-shot)~\citep{Eval_ARC}, Arc-Challenge (0-shot)~\citep{Eval_ARC}, HellaSwag (0-shot)~\citep{Eval_hellaswag}, PIQA (0-shot)~\citep{Eval_piqa}, and MMLU (5-shot)~\citep{hendrycks2020measuring}. For evaluation, we use a prompt likelihood-based scoring approach, computing probabilities for each candidate answer, and present the mean accuracy as the metric.
\paragraph{Coding related generation tasks:} Pass@1 scores are provided for MBPP (3-shot)~\citep{Eval_MBPP} and HumanEval (0-shot)~\citep{Eval_HumanEval}, serving as indicators for the model’s capacity to synthesize Python programs.

\autoref{tab:evals} presents a comparison among three models, each trained on the {\textsc BLT}-1T dataset: a model based on the Llama 3 tokenizer using bpe,\footnote{We opt for the Llama 3 tokenizer, featuring a 128k vocabulary, as it outperforms the 32k vocabulary version found in Llama 2.} as well as two {\textsc BLT} model variants. The first variant adopts a space-patching approach ({\textsc BLT}-Space), while the second leverages an entropy-based patching strategy ({\textsc BLT}-Entropy), employing an approximate monotonicity condition and resetting the model’s context at new line characters, as outlined in~\autoref{section:entropy-context}. Training for all three configurations utilizes a comparable computational budget (flops). Notably, for the {\textsc BLT}-Entropy model, we implement a modification at inference by lowering the entropy threshold from 0.6 to 0.1. This adjustment yields enhanced performance on downstream tasks, although it requires a greater number of inference steps.

The {\textsc BLT}-Entropy model achieves superior performance compared to the Llama 3 model on four of the seven evaluated tasks, despite both being trained on an equal quantity of bytes. This enhancement can likely be attributed to two factors: (1) more efficient utilization of training compute enabled by dynamic patching, and (2) the model’s direct approach to capturing byte-level data rather than relying on token-based representations.

Conversely, {\textsc BLT}-Space delivers inferior results relative to the Llama 3 tokenizer across all except one task; however, it substantially decreases inference flops owing to its greater average patch size of 6 bytes. By contrast, both the bpe and entropy-patching models demonstrate similar average patch sizes of around 4.5 bytes within the training dataset. Given the same training budget, the model with the larger patch size is able to process 30\% more data than the other two approaches, potentially moving {\textsc BLT} further from the compute-optimal regime.

\subsection{Patches Scale Better Than Tokens} In {\textsc BLT} models, both the size of the model and the size of each patch can be increased at the same time, all while preserving the training and inference computation budget and keeping the volume of training data unchanged. This flexibility in enlarging patch size at will is distinctive to patch-based architectures, liberating them from the efficiency limitations imposed by token-based models with a fixed vocabulary, as elaborated in Section~\ref{section:bpe-patch}. Employing larger patches conserves computational resources, making it possible to allocate more computation towards scaling the global latent transformer, which operates with reduced frequency as patch size grows.

In order to evaluate the proposition that, at fixed inference cost, increasing model size and patch dimensions—while reducing step count—might yield superior performance relative to smaller models with more processing steps, we conduct a scaling experiment. Using starting model sizes of 400 million and 3.6 billion parameters, both employing the Llama 2 tokenizer, we identify models with comparable computational footprint under the Llama 3 tokenizer as well as {\textsc BLT}-Entropy variants configured with mean patch sizes of 6 and 8 bytes over the same training data mixture (see~\autoref{tab:fixed-inf-params} for further model specifications). For the models with patch size 8, three encoder layers are utilized in lieu of one. All models are trained across a range of allocated training FLOP budgets.

As illustrated in \autoref{fig:fixed_inference_scaling}, {\textsc BLT} models display superior scaling behavior compared to tokenization-based models across both categories of inference flops. Although BPE models demonstrate stronger performance with limited training resources, they are quickly outpaced by {\textsc BLT} models once the training budget extends only slightly beyond the compute-optimal point. From a practical standpoint, investing more in initial pretraining often leads to higher-performing models under a constrained inference budget. A prominent illustration of this is the 8B model class, such as Llama 3.1, which was trained on data volumes two orders of magnitude greater than the compute-optimal amount for that model size.

The crossover threshold at which {\textsc BLT} begins to surpass token-based architectures has moved marginally closer to the point of compute optimality as we examine models belonging to the higher FLOP category, shifting from approximately 3x to 2.5x the optimal compute allocation. In addition, the model employing a larger patch size of 8 demonstrates a more pronounced scaling curve within the higher FLOP class, allowing it to surpass competing models at an earlier stage. As elaborated in Section~\ref{section:parameter-matched-scaling}, it appears that, at greater model sizes, configurations with larger patch sizes approach the performance of BPE models. We partially ascribe this to the fact that the proportion of total FLOPs consumed by the byte-level Encoder and Decoder modules diminishes, as these modules scale less rapidly in comparison to the Latent Transformer. When we increase the total parameter count by a factor of 20—from 400M up to 8B—this only results in approximately a twofold increase in {\textsc BLT}'s local model parameters. This consideration is critical because augmenting the patch size primarily impacts the FLOPs associated with the patch Latent Transformer, leaving the byte-level components largely unaffected. Notably, this accounts for the observation that the {\textsc BLT}-Entropy ps=8 model’s relative size, compared to the Llama 2, increased from 1.6x to 1.7x when scaling up to the larger model.

In conclusion, our investigation into patch-length scaling reveals that the {\textsc BLT} patch-based framework exhibits enhanced scaling behavior when both patch size and model capacity are enlarged together. Notably, this positive scaling persists, and in some cases intensifies, as the models are further scaled up.

\section{Byte Modeling Improves Robustness} We further evaluate the robustness of {\textsc BLT} in comparison with token-based models that do not incorporate direct byte-level signals, and outline a method for transforming pretrained token-based architectures to operate at the byte level.
\label{sec:robustness}

\subsection{Character-Level Tasks}

Initial motivations for the development of byte-level models included leveraging their resilience to noise introduced at the byte level within input data, as well as their sensitivity to the internal composition of tokens—challenges that remain for existing models based on tokenizers. To investigate these attributes, we conduct supplementary assessments using benchmarks designed to test both the robustness against input perturbations and the recognition of character-level information, spanning English, multilingual data, digits, and phonetic symbols. The outcomes of these evaluations are detailed in Table~\ref{tab:char_tasks}.

\paragraph{Noisy Data} To assess the robustness of {\textsc BLT} relative to tokenizer-based models, we generate noisy variants of the benchmark classification datasets discussed in Section~\ref{section:evals}. Five character-level noising methods are implemented to perturb the input text as follows: (a) \textit{AntSpeak}, where text is transformed entirely into uppercase and characters are separated by spaces; (b) \textit{Drop}, in which 10\% of characters are randomly eliminated; (c) \textit{RandomCase}, randomly converts half of all characters to uppercase and the remaining half to lowercase; (d) \textit{Repeat}, duplicates 20\% of the characters up to four times; and (e) \textit{UpperCase}, which capitalizes the entire text. For evaluation, each strategy is individually applied to the prompt, the completion, or both, treating each configuration as a distinct task, and results are averaged. As presented in Table~\ref{tab:char_tasks}, experiments on the noised HellaSwag~\citep{Eval_hellaswag} dataset reveal that {\textsc BLT} consistently demonstrates superior robustness in comparison to tokenizer-based models, achieving a mean performance boost of 8 points over the model trained with identical data, and surpassing the Llama 3.1 model, which benefits from training on a significantly larger corpus.

\paragraph{Phonology - Grapheme-to-Phoneme (G2P)} We evaluate how effectively {\textsc BLT} converts sequences of graphemes—that is, the written characters of a word—into their corresponding phoneme transcriptions. Table \ref{tab:char_tasks} displays the results for the G2P task conducted under a 5-shot learning regime with Phonology Bench~\citep{suvarna2024phonologybench}. Our analysis indicates that {\textsc BLT} achieves superior performance compared to the baseline Llama 3 1T tokenizer-based model on this task.

\paragraph{CUTE}
To gauge the extent of character-level understanding, we test {\textsc BLT} using the CUTE benchmark~\citep{edman2024cute}, which contains a diverse suite of tasks grouped into three main categories: composition comprehension, orthographic similarity discernment, and sequence manipulation capability. CUTE is notably challenging for most models that rely on tokenizers, as while these models are aware of how their tokens are spelled, they often fall short when manipulating textual sequences based on that knowledge. As presented in Table~\ref{tab:char_tasks}, {\textsc BLT}-Entropy surpasses both BPE Llama 3 variants by a margin of over 25 points on this suite. Notably, {\textsc BLT} excels at tasks involving manipulation at the character level, obtaining a striking 99.9\% on both spelling-oriented challenges. These substantial performance gains, even with {\textsc BLT} being trained on just one-sixteenth the data volume used by Llama 3.1, suggest that learning character-level structures proves particularly difficult for BPE-based architectures. Figure~\ref{fig:cute_examples} showcases cases where the Llama 3 tokenizer encounters difficulties, whereas our {\textsc BLT} demonstrates robust performance. The only areas where BPE outperforms are word deletion and insertion tasks; such types of manipulations may not be as natural for a byte-level architecture, though the performance gap is not insurmountable. In fact, synthesizing words from characters may prove to be simpler than decomposing words into finer subunits. Throughout all assessments, we adopt an identical evaluation protocol and utilize the original Huggingface prompts. It is plausible that BPE models could achieve higher performance with enhanced prompt engineering.

\paragraph{Low Resource Machine Translation} We assess the performance of {\textsc BLT} on translation tasks involving six major language families and twenty-one low-resource languages using various scripts, drawing on the FLORES-101 benchmark~\citep{goyal2022flores}. The SentencePiece BLEU scores are presented in Table~\ref{tab:flores}. The findings indicate that {\textsc BLT} surpasses the model utilizing the Llama 3 tokenizer, offering an overall improvement of 2 BLEU points for translations into English and a 0.5-point gain when translating from English. For prominent language pairs, {\textsc BLT}'s results are on par with, or marginally surpass, those achieved by Llama 3. Notably, for numerous low-resource language pairs, {\textsc BLT} achieves superior performance compared to Llama 3, highlighting the capability of byte-level modeling to effectively generalize across rare and complex byte sequences.

\subsection{Training {\textsc BLT} from Llama 3}
\label{sec:distillation}
We investigate a methodology where {\textsc BLT} models capitalize on established pre-trained tokenizer-based models to facilitate more efficient and rapid convergence during training. This is implemented by initializing the global transformer parameters of {\textsc BLT} directly from those of a pre-trained Llama 3.1 model. In the subsequent training phase, we fine-tune the global transformer using a learning rate that is one-tenth of what is used for the local encoder and local decoder, and for the number of steps optimized for Llama 3. Table~\ref{tab:distill} presents the results in comparison to a baseline {\textsc BLT}. The outcomes demonstrate that {\textsc BLT} initialized from Llama 3.1 yields substantial gains over both the Llama 3 and the standard {\textsc BLT} baselines, under equivalent FLOP budgets. Furthermore, in contrast with our {\textsc BLT}-Entropy variant (reported in Table~\ref{tab:evals}), which received training on a much larger corpus (1T tokens), {\textsc BLT} transferred from Llama 3.1 still achieves stronger performance on the MMLU benchmark. This observation underscores the approach’s potential for dramatically reducing training FLOPs while attaining improved effectiveness.

This configuration can alternatively be interpreted as adapting tokenizer-dependent models into tokenizer-free architectures, thereby transforming the pre-trained LLaMA 3.1 into a {\textsc BLT} model. For a thorough comparison, we present results for the original LLaMA 3.1 model, trained with 15T tokens, in Table \ref{tab:distill}, and assess it alongside the {\textsc BLT} obtained from LLaMA 3. Our approach shows a minor reduction in performance on MMLU and HumanEval, whereas more pronounced decreases are observed for the remaining tasks. These findings indicate that additional research is required to more effectively utilize the pre-trained model, with particular emphasis on optimizing data compositions and other relevant hyperparameters to enhance overall outcomes.

\section{Ablations and Discussion}
\label{section:ablations}
Here, we present ablation studies to motivate our design decisions for {\textsc BLT}, including the selection of the patching strategy and tuning of hyperparameters used in training the 8B parameter {\textsc BLT} model on the {\textsc BLT}-1T dataset.

\paragraph{Entropy Model Hyper-parameters} In order to analyze how adjustments in entropy model capacity and context window size influence scaling behavior, we experiment with byte-level entropy transformer models that range from 1m to 100m parameters, employing context windows sized from 64 up to 512. To visualize the results, we generate curves of bits-per-byte (bpb) against training flops, utilizing our $400m$ and $1b$ {\textsc BLT} models trained on the Llama-2 corpus. These curves are displayed in \autoref{fig:entropy_ablation}. Our findings indicate that scaling along both axes—model parameter count and context window length—enhances performance, though performance improvements plateau once surpassing 50m model parameters.

\paragraph{Types of Patching}

We conduct an ablation study of the four patching strategies outlined in Section~\ref{section:patching}, namely: 1) Strided Patching employing strides of 4 and 6, 2) Whitespace-based Patching, 3) Patching utilizing BPE Tokenization aligned with the Llama 3 tokenizer, and 4) Entropy-driven patching leveraging a compact byte-level LLM.

Although dynamic patching shortens sequence lengths, we standardize the context length across all patching methods by carefully controlling sequence length. This ensures that, on average, each model processes an equivalent number of bytes per sequence during both training and inference, thereby eliminating confounding influences arising from differing context sizes. As shown in Figure~\ref{fig:scaling-compute-opt}, these ablation studies reveal that every alternative patching strategy surpasses static patching in performance, with space patching performing nearly as well as dynamic entropy-based patching.

\autoref{tab:evals_ablation} displays the results of benchmark experiments for {\textsc BLT} models, contrasting the performance of tokenizer-based, space patching, and entropy-based patching approaches. These models were all trained on the Llama 2 dataset for the empirically determined optimal number of steps~\citep{dubey2024llama}. Despite the relative simplicity of space patching—which does not require real-time evaluation of an entropy model during training—we observe that the improvements associated with entropy-based patching in the context of scaling (see Section~\ref{section:scaling}) also manifest on downstream benchmark evaluations.\footnote{For space patching, results correspond to earlier experiments lacking cross-attention, though similar performance patterns persist even when cross-attention is applied.}

\paragraph{Cross-Attention} 
In~\autoref{tab:crossattn_ablations_1b}, we conduct an ablation study examining the impact of integrating cross-attention modules at different locations within both the encoder and decoder of {\textsc BLT}. Specifically, for the encoder side cross-attention, we experiment with three strategies for initializing the queries: 1) employing a single learned embedding that is identical across all global states, 2) using a hash-based embedding derived from the patch's byte sequence, and 3) applying a pooling operation on the encoder’s hidden states corresponding to the patch bytes at the specific encoder layer in question.

Our results indicate that incorporating cross-attention within the \textit{decoder} yields the greatest benefit. Employing cross-attention in the encoder offers a modest gain, which arises only when queries are initialized through pooling. Moreover, cross-attention proves especially advantageous on the Common-Crawl dataset and is even more beneficial when utilizing larger patch sizes.

\paragraph{n-gram Hash Embeddings} 

We evaluate the impact of varying the n-gram hash embedding vocabulary sizes (0, 100K, 200K, and 400K), with outcomes shown in Table~\ref{tab:ngram_ablations_1b}. Our analysis demonstrates that incorporating hash embeddings consistently enhances results across all tested domains, with the most noticeable benefits seen on Wikipedia and Github, where there is a 0.04 bpb improvement compared to 0.01 bpb after 15k training steps at the 8B scale. At the same 8B scale, reducing the hash size from 500K to 300K only led to a minor change in performance of around 0.001 bpb after 15k steps. These findings suggest that hash embeddings play a crucial role in narrowing the gap between {\textsc BLT} and traditional tokenizer-based models, though increasing the hash vocabulary beyond 300K yields minimal additional benefits. Furthermore, the observed improvements from hash embeddings tend to supplement those provided by cross-attention mechanisms, as each confers gains on distinct datasets.

\paragraph{Local Model Hyperparameters} Table \ref{tab:local_layers} presents an ablation study investigating different choices for the layer counts in both the local encoder and decoder. Notably, {\textsc BLT} demonstrates strong performance when using hash n-gram embeddings in combination with a highly simplified encoder—utilizing only a single layer—while benefitting from a more complex, deeper decoder architecture.

\section{Related Work}
\label{section:related_work}

\textbf{Character-Level RNNs:} Character language modeling has maintained its relevance since the inception of neural language models~\citep{sutskever2011generating,mikolov2012subword,graves2013generating}, largely due to its inherent ability to naturally handle out-of-vocabulary terms, thus obviating the need for back-off strategies. \cite{kim2016character} introduced an approach that exclusively utilizes character-level information on the input side, applying convolutional and highway networks that subsequently inform LSTM-based RNNs; this architecture demonstrated competitive performance relative to leading RNN-based language models for English and achieved superior results for morphologically complex languages—highlighting another notable benefit of character-level LLMs. In another study, \cite{kenter2018byte} presented byte-level LSTM models for machine comprehension, yielding better outcomes than their word-level counterparts particularly for morphologically rich languages such as Turkish and Russian. Continuing in this vein, \cite{zhang2015character} leveraged character-level convolutional architectures for classification problems, achieving higher accuracy on specific tasks compared to models relying on word-level representations. \citet{chung2019hierarchical} employed hierarchical LSTM structures incorporating boundary detectors at each layer, enabling the discovery of latent hierarchical structures within text, thereby enhancing the effectiveness of character-level language modeling. In a distinct line of work, ByteNet~\cite{kalchbrenner2016neural} utilized CNN layers directly on character sequences instead of leveraging attention mechanisms for tasks in machine translation.

\textbf{Character-Level Transformers:} The integration of attention mechanisms in transformer architectures~\citep{vaswani2017attention} in combination with subword tokenization methods~\citep{sennrich-etal-2016-neural} has notably enhanced the capabilities of neural models in language modeling and other standard tasks. Nonetheless, employing word and subword units enforces an implicit inductive bias concerning the granularity at which models are expected to process input. To bring together the advantages offered by transformer frameworks and the early encouraging outcomes in character-level language modeling, \citet{al2019character} explored the use of very deep transformer networks. By incorporating auxiliary loss functions, they trained transformer models on character data that surpassed the previous LSTM-based character language models in performance. Despite these improvements, their results still exhibited a notable discrepancy when compared to word-level LLMs. Similarly, GPT-2~\citep{radfordgpt2} reported that on large-scale datasets, such as the 1 Billion Word Benchmark, byte-level language models failed to match the effectiveness of their word-level counterparts.

Although \cite{Choe2019BridgingTG} showed that transformer-based byte-level LLMs can surpass subword-level counterparts with similar model sizes, they noted a substantial increase in computation requirements and training time. Likewise, \cite{el2020characterbert} introduced CharFormer, a BERT-style architecture that generates word representations through convolutions over character embeddings, demonstrating its efficacy in medical text analysis, albeit at a much greater computational expense. In another approach, \cite{clark2022canine} proposed CANINE, an encoder-only model comprising 150M parameters, which processes raw character sequences. CANINE leverages a deep transformer structure—paralleling the design of our global model—and incorporates local transformers together with strided convolutions for character downsampling. This setup enables CANINE to outperform a parameter-matched token-level model (mBERT) on a variety of multilingual downstream tasks. ByT5~\citep{xue2022byt5} investigated byte-level encoder-decoder models that eschew any kind of patching mechanism. While ByT5 displayed enhanced robustness to input noise and achieved results comparable to those of token-based models using only a quarter of the data, the absence of input patching led to prohibitively expensive attention computations across every byte, making the model computationally inefficient. Modeling raw bytes, as opposed to subword units, significantly lengthens the sequence input, compounding the difficulty of scaling such byte-level models efficiently. Recently, \cite{wang2024mambabyte} leveraged the Mamba Architecture~\citep{gu2023mamba}, which is capable of sustaining a fixed-size memory over exceptionally long contexts, to train a byte-level Mamba model without patching. At the 350M parameter scale, their approach surpasses byte-level transformers—within matched FLOPs constraints—when evaluated by bits-per-byte across a variety of datasets.

\textbf{Patching-based approaches:} Leveraging patching methods can substantially reduce the typically high computational cost (in terms of flops) incurred by byte-level LLMs, while also potentially preserving their performance. Several studies have provided early evidence of success with these techniques at smaller model and dataset scales. For instance, \cite{nawrot-etal-2022-hierarchical} utilize a static approach to downsampling and upsampling through patching, resulting in the hourglass transformer architecture, which surpasses alternative byte-level models at the 150M parameter scale. Subsequently, \cite{nawrot-etal-2023-efficient} build upon these ideas by introducing dynamic patching mechanisms, such as an end-to-end trainable boundary predictor, a boundary predictor trained with explicit supervision from certain tokenizers, and an entropy-driven patching strategy akin to {\textsc BLT}. Their findings indicate that such approaches can outperform contemporaneous vanilla transformer models on language modeling benchmarks, specifically at the 40M parameter range and with training data comprising approximately 400M tokens. In a different line of work, \cite{lester2024training} explore the use of arithmetic coding to compress sequences during training, achieving compression ratios superior to those obtained via BPE. By applying an equal-info windows methodology, their models exceed the performance of byte-level baselines on language modeling tasks, though they still fall short compared to subword-level baselines.

Our research is heavily influenced by and bears the closest resemblance to MegaByte~\citep{yu2023megabyte}, a decoder-only causal LLM architecture that employs a fixed, static patching scheme. In this setup, bytes are converted into patches by concatenating representations, and on the decoder side, a local model reverts these patches to their original byte format. MegaByte was shown to achieve performance on par with tokenizer-based architectures at a 1B parameter scale when evaluated on a dataset comprising approximately 400B bytes. In all our experimental analyses, we include comprehensive ablations of MegaByte, observing that its static patching approach falls short of the best compute-optimized, tokenizer-based models in scenarios where compute is strictly regulated. We illustrate how {\textsc BLT} addresses this discrepancy. Similar conclusions regarding MegaByte's performance are drawn by \cite{slagle2024spacebyte}, who propose expanding static patching to include segmentation on whitespace and similar space-like bytes, along with incorporating a local encoder model. They report gains over tokenizer-based transformer models under compute-constrained settings within specific domains like Github and arXiv, again at the 1B parameter level. Our own experiments using this model further indicate that advancing the architectural design is crucial for byte-level models to scale effectively and rival leading-edge token-based models such as Llama 3.

\section{Limitations and Future Work}

In this study, we adopt the number of training steps identified as optimal for Llama~3~\citep{dubey2024llama} to guide our architectural decisions. Nonetheless, it should be noted that these scaling laws were originally developed for BPE-level transformer models and may not yield ideal ratios between data and model size for {\textsc BLT}. Determining scaling laws specifically for {\textsc BLT}—which could potentially reveal more advantageous scaling behavior for our model—remains an open avenue for future investigation. Furthermore, the majority of our experiments were performed on models up to 1B parameters, and it is plausible that optimal architectural settings might differ at larger model scales, such as 8B parameters or higher, thereby offering the potential for enhanced performance in larger regimes.

Current transformer codebases and libraries are optimized primarily for tokenizer-based transformer models. Although our work includes theoretical comparisons using matched FLOPs and incorporates some efficient implementations—such as FlexAttention—to manage layers that differ from standard transformer configurations, our practical implementations might still lag behind tokenizer-based approaches regarding actual wall-clock performance, and could see improvement through additional optimization efforts.

Whereas {\textsc BLT} employs an independently trained entropy model for the patching process, integrating the learning of this patching model into an end-to-end training paradigm represents a promising avenue for future exploration. Section \ref{sec:distillation} details preliminary experiments that provide evidence for the feasibility of “byte-ifying” tokenizer-based architectures, including Llama 3, which have been trained on datasets exceeding 10T tokens. This is achieved by initializing and keeping the global transformer parameters fixed using their pre-trained weights. Continued research along these lines holds the potential not only to preserve the advantages associated with bytefying, but also to enhance the performance of such models—potentially even surpassing that of the original tokenizer-based models—without necessitating full retraining from the ground up.

\section{Conclusion}
\label{section:conclusion}

In this work, we introduce the Byte Latent Transformer (\textbf{BLT}), an innovative model architecture that challenges the prevailing reliance on predetermined vocabulary tokenization in large language models. BLT incorporates a flexible, learnable mechanism for assembling bytes into variable-sized patches, allowing the allocation of computational resources to be dynamically aligned with data complexity. This design yields substantial gains in both computational efficiency and model robustness. Through a comprehensive scaling investigation, we establish that BLT models attain parity with tokenization-based counterparts such as Llama 3 at scales of up to 8B parameters and 4T bytes, and can achieve inference FLOP reductions of up to 50\% with only marginal decreases in evaluation scores. Moreover, BLT facilitates a novel scaling axis: it enables the expansion of both model and patch dimensions under a fixed inference cost, offering practical advantages in compute-constrained environments. By directly processing raw byte sequences, BLT enhances the system’s capacity to cope with the rarer aspects of input data, yielding notable robustness against noise as well as a finer-grained modeling of sub-word elements. Collectively, these findings suggest that BLT is a strong candidate for replacing traditional tokenization schemes, paving the way for highly efficient, scalable, and resilient language models.

\end{document}